% Options for packages loaded elsewhere
\PassOptionsToPackage{unicode}{hyperref}
\PassOptionsToPackage{hyphens}{url}
\documentclass[
  11pt,
  a4paper,
]{article}
\usepackage{xcolor}
\usepackage[margin=18mm]{geometry}
\usepackage{amsmath,amssymb}
\setcounter{secnumdepth}{-\maxdimen} % remove section numbering
\usepackage{iftex}
\ifPDFTeX
  \usepackage[T1]{fontenc}
  \usepackage[utf8]{inputenc}
  \usepackage{textcomp} % provide euro and other symbols
\else % if luatex or xetex
  \usepackage{unicode-math} % this also loads fontspec
  \defaultfontfeatures{Scale=MatchLowercase}
  \defaultfontfeatures[\rmfamily]{Ligatures=TeX,Scale=1}
\fi
\ifPDFTeX\else
  % xetex/luatex font selection
    \setmainfont[]{Noto Serif}
    \setsansfont[]{Noto Sans}
    \setmonofont[]{Noto Sans Mono}
  \setmathfont[]{Noto Sans Math}
\fi
% Use upquote if available, for straight quotes in verbatim environments
\IfFileExists{upquote.sty}{\usepackage{upquote}}{}
\IfFileExists{microtype.sty}{% use microtype if available
  \usepackage[]{microtype}
  \UseMicrotypeSet[protrusion]{basicmath} % disable protrusion for tt fonts
}{}
\makeatletter
\@ifundefined{KOMAClassName}{% if non-KOMA class
  \IfFileExists{parskip.sty}{%
    \usepackage{parskip}
  }{% else
    \setlength{\parindent}{0pt}
    \setlength{\parskip}{6pt plus 2pt minus 1pt}}
}{% if KOMA class
  \KOMAoptions{parskip=half}}
\makeatother
\ifLuaTeX
\usepackage[bidi=basic]{babel}
\else
\usepackage[bidi=default]{babel}
\fi
\babelprovide[main,import]{russian}
\ifPDFTeX
\else
\babelfont{rm}[]{Noto Serif}
\fi
% get rid of language-specific shorthands (see #6817):
\let\LanguageShortHands\languageshorthands
\def\languageshorthands#1{}
\setlength{\emergencystretch}{3em} % prevent overfull lines
\providecommand{\tightlist}{%
  \setlength{\itemsep}{0pt}\setlength{\parskip}{0pt}}
\usepackage{bookmark}
\IfFileExists{xurl.sty}{\usepackage{xurl}}{} % add URL line breaks if available
\urlstyle{same}
\hypersetup{
  pdftitle={Билеты по дискретным случайным процессам},
  pdflang={ru-RU},
  hidelinks,
  pdfcreator={LaTeX via pandoc}}

\title{Билеты по дискретным случайным процессам}
\author{}
\date{}

\begin{document}
\maketitle

\section{Билет 24. Корреляционная функция. Достаточные условия
эргодичности
последовательности}\label{ux431ux438ux43bux435ux442-24.-ux43aux43eux440ux440ux435ux43bux44fux446ux438ux43eux43dux43dux430ux44f-ux444ux443ux43dux43aux446ux438ux44f.-ux434ux43eux441ux442ux430ux442ux43eux447ux43dux44bux435-ux443ux441ux43bux43eux432ux438ux44f-ux44dux440ux433ux43eux434ux438ux447ux43dux43eux441ux442ux438-ux43fux43eux441ux43bux435ux434ux43eux432ux430ux442ux435ux43bux44cux43dux43eux441ux442ux438}

Источники: Ширяев 2, гл. V-VI; рукописные лекции ДСП, лекции 9-10;
Сердобольская, разделы о ковариационной функции.

\subsection{1. Короткий
ответ}\label{ux43aux43eux440ux43eux442ux43aux438ux439-ux43eux442ux432ux435ux442}

Для случайной последовательности со вторыми моментами ковариационная
функция описывает линейную зависимость значений в разные моменты
времени. Если последовательность стационарна в широком смысле и

\[
\mathbb E X_n=m,
\]

то ее ковариационная функция зависит только от лага:

\[
R(k)=\operatorname{Cov}(X_{n+k},X_n)
=
\mathbb E[(X_{n+k}-m)\overline{(X_n-m)}].
\]

Если \(R(0)>0\), нормированная корреляционная функция:

\[
\rho(k)=\frac{R(k)}{R(0)}.
\]

Эргодичность по среднему в широком смысле означает

\[
\overline X_N=\frac1N\sum_{n=0}^{N-1}X_n
\xrightarrow[L^2]{N\to\infty}
m.
\]

Главная формула:

\[
\mathbb E|\overline X_N-m|^2
=
\frac1{N^2}\sum_{i=0}^{N-1}\sum_{j=0}^{N-1}R(i-j)
=
\frac1N\sum_{|k|<N}\left(1-\frac{|k|}{N}\right)R(k).
\]

Поэтому достаточное условие эргодичности:

\[
\frac1N\sum_{|k|<N}\left(1-\frac{|k|}{N}\right)R(k)\to0.
\]

Более простые достаточные условия:

\[
R(k)\to0
\]

или

\[
\sum_{k\in\mathbb Z}|R(k)|<\infty.
\]

\subsection{2. Полный
ответ}\label{ux43fux43eux43bux43dux44bux439-ux43eux442ux432ux435ux442}

\subsubsection{2.1. Введение и связь с предыдущими
темами}\label{ux432ux432ux435ux434ux435ux43dux438ux435-ux438-ux441ux432ux44fux437ux44c-ux441-ux43fux440ux435ux434ux44bux434ux443ux449ux438ux43cux438-ux442ux435ux43cux430ux43cux438}

Этот билет продолжает Билет 17 о стационарности и эргодичности. Там были
даны определения стационарности в узком и широком смыслах. Здесь нужно
разобрать, как эргодичность среднего проверяется через корреляционную,
точнее ковариационную, функцию.

\subsubsection{2.2. Определение ковариации и стационарность в широком
смысле}\label{ux43eux43fux440ux435ux434ux435ux43bux435ux43dux438ux435-ux43aux43eux432ux430ux440ux438ux430ux446ux438ux438-ux438-ux441ux442ux430ux446ux438ux43eux43dux430ux440ux43dux43eux441ux442ux44c-ux432-ux448ux438ux440ux43eux43aux43eux43c-ux441ux43cux44bux441ux43bux435}

Пусть

\[
X=(X_n)_{n\in\mathbb Z}
\]

\begin{itemize}
\tightlist
\item
  случайная последовательность с конечными вторыми моментами:
\end{itemize}

\[
\mathbb E|X_n|^2<\infty.
\]

Для двух случайных величин \(X_n\) и \(X_m\) ковариация равна

\[
\operatorname{Cov}(X_n,X_m)
=
\mathbb E[(X_n-\mathbb EX_n)\overline{(X_m-\mathbb EX_m)}].
\]

Для действительных величин комплексное сопряжение можно не писать.

Если последовательность стационарна в широком смысле, то:

\[
\mathbb EX_n=m
\]

не зависит от \(n\), а ковариация зависит только от разности индексов:

\[
\operatorname{Cov}(X_n,X_m)=R(n-m).
\]

Обычно фиксируют

\[
R(k)=\operatorname{Cov}(X_k,X_0).
\]

Если процесс центрирован, то есть \(m=0\), то

\[
R(k)=\mathbb E X_k\overline{X_0}.
\]

В рукописных лекциях это записывается как скалярное произведение в
гильбертовом пространстве \(L_2\):

\[
R(k)=(X_k,X_0)_{L_2}.
\]

\subsubsection{2.3. Различие между ковариационной и корреляционной
функциями}\label{ux440ux430ux437ux43bux438ux447ux438ux435-ux43cux435ux436ux434ux443-ux43aux43eux432ux430ux440ux438ux430ux446ux438ux43eux43dux43dux43eux439-ux438-ux43aux43eux440ux440ux435ux43bux44fux446ux438ux43eux43dux43dux43eux439-ux444ux443ux43dux43aux446ux438ux44fux43cux438}

Важно различать две функции:

\begin{enumerate}
\def\labelenumi{\arabic{enumi}.}
\tightlist
\item
  ковариационная функция:
\end{enumerate}

\[
R(k)=\operatorname{Cov}(X_k,X_0);
\]

\begin{enumerate}
\def\labelenumi{\arabic{enumi}.}
\setcounter{enumi}{1}
\tightlist
\item
  корреляционная функция в нормированном смысле:
\end{enumerate}

\[
\rho(k)=\frac{R(k)}{R(0)}
\]

при \(R(0)>0\).

В разных курсах слово ``корреляционная'' иногда используют для \(R(k)\),
а иногда для \(\rho(k)\). На экзамене лучше сразу сказать: ``я буду
называть \(R(k)\) ковариационной функцией, а нормированную величину
\(\rho(k)\) - корреляционной''.

\subsubsection{2.4. Основные свойства ковариационной
функции}\label{ux43eux441ux43dux43eux432ux43dux44bux435-ux441ux432ux43eux439ux441ux442ux432ux430-ux43aux43eux432ux430ux440ux438ux430ux446ux438ux43eux43dux43dux43eux439-ux444ux443ux43dux43aux446ux438ux438}

Ковариационная функция имеет несколько базовых свойств.

Во-первых,

\[
R(0)=\operatorname{Var}X_0\ge0.
\]

Если

\[
R(0)=0,
\]

то

\[
X_n=m
\]

почти наверное для всех \(n\), то есть процесс тривиален в смысле
второго порядка.

Во-вторых, для комплексного случая

\[
R(-k)=\overline{R(k)}.
\]

Для действительного процесса:

\[
R(-k)=R(k).
\]

В-третьих, из неравенства Коши-Буняковского:

\[
|R(k)|
=
|\operatorname{Cov}(X_k,X_0)|
\le
\sqrt{\operatorname{Var}X_k\,\operatorname{Var}X_0}
=
R(0).
\]

В-четвертых, \(R\) положительно определена. Для любых комплексных чисел
\(a_1,\ldots,a_r\) и целых индексов \(n_1,\ldots,n_r\):

\[
\sum_{\alpha=1}^r\sum_{\beta=1}^r
a_\alpha \overline{a_\beta}R(n_\alpha-n_\beta)
\ge0.
\]

Это просто дисперсия линейной комбинации:

\[
\mathbb E\left|
\sum_{\alpha=1}^r a_\alpha (X_{n_\alpha}-m)
\right|^2
\ge0.
\]

\subsubsection{2.5. Понятие эргодичности по среднему в широком
смысле}\label{ux43fux43eux43dux44fux442ux438ux435-ux44dux440ux433ux43eux434ux438ux447ux43dux43eux441ux442ux438-ux43fux43e-ux441ux440ux435ux434ux43dux435ux43cux443-ux432-ux448ux438ux440ux43eux43aux43eux43c-ux441ux43cux44bux441ux43bux435}

Теперь эргодичность. В широком смысле чаще всего рассматривают
эргодичность по среднему: можно ли среднее по одной длинной реализации
заменить математическим ожиданием?

Определим временное среднее:

\[
\overline X_N
=
\frac1N\sum_{n=0}^{N-1}X_n.
\]

Последовательность называется эргодичной по среднему в широком смысле,
если

\[
\overline X_N
\xrightarrow[L^2]{N\to\infty}
m.
\]

То есть

\[
\mathbb E|\overline X_N-m|^2\to0.
\]

Так как

\[
\mathbb E\overline X_N=m,
\]

это равносильно

\[
\operatorname{Var}(\overline X_N)\to0.
\]

\subsubsection{2.6. Вывод формулы для дисперсии временного
среднего}\label{ux432ux44bux432ux43eux434-ux444ux43eux440ux43cux443ux43bux44b-ux434ux43bux44f-ux434ux438ux441ux43fux435ux440ux441ux438ux438-ux432ux440ux435ux43cux435ux43dux43dux43eux433ux43e-ux441ux440ux435ux434ux43dux435ux433ux43e}

Вычислим эту дисперсию. Для центрированной последовательности
\(Y_n=X_n-m\):

\[
\overline X_N-m
=
\frac1N\sum_{n=0}^{N-1}Y_n.
\]

Тогда

\[
\mathbb E|\overline X_N-m|^2
=
\frac1{N^2}
\sum_{i=0}^{N-1}\sum_{j=0}^{N-1}
\mathbb E Y_i\overline{Y_j}.
\]

По стационарности в широком смысле:

\[
\mathbb E Y_i\overline{Y_j}=R(i-j).
\]

Значит

\[
\mathbb E|\overline X_N-m|^2
=
\frac1{N^2}
\sum_{i=0}^{N-1}\sum_{j=0}^{N-1}R(i-j).
\]

Если собрать пары с одинаковым лагом \(k=i-j\), то число таких пара
равно \(N-|k|\). Поэтому

\[
\mathbb E|\overline X_N-m|^2
=
\frac1N
\sum_{|k|<N}
\left(1-\frac{|k|}{N}\right)R(k).
\]

Для действительной последовательности это можно записать как

\[
\operatorname{Var}(\overline X_N)
=
\frac{R(0)}{N}
+
\frac{2}{N}
\sum_{k=1}^{N-1}
\left(1-\frac{k}{N}\right)R(k).
\]

\subsubsection{2.7. Основной критерий эргодичности и его спектральная
интерпретация}\label{ux43eux441ux43dux43eux432ux43dux43eux439-ux43aux440ux438ux442ux435ux440ux438ux439-ux44dux440ux433ux43eux434ux438ux447ux43dux43eux441ux442ux438-ux438-ux435ux433ux43e-ux441ux43fux435ux43aux442ux440ux430ux43bux44cux43dux430ux44f-ux438ux43dux442ux435ux440ux43fux440ux435ux442ux430ux446ux438ux44f}

Отсюда сразу получаем главный критерий второго порядка:

\[
\overline X_N\to m\ \text{в }L^2
\quad\Longleftrightarrow\quad
\frac1{N^2}
\sum_{i=0}^{N-1}\sum_{j=0}^{N-1}R(i-j)\to0.
\]

Для центрированной стационарной \(L^2\)-последовательности это можно
формулировать через среднее Чезаро ковариаций:

\[
\frac1N\sum_{k=0}^{N-1}R(k)\to0
\]

как необходимое и достаточное условие сходимости средних арифметических
в среднеквадратическом смысле. В спектральном языке это эквивалентно
отсутствию атома спектральной меры в нуле:

\[
F(\{0\})=0.
\]

Но для билета достаточно знать простые достаточные условия.

\subsubsection{2.8. Достаточные условия
эргодичности}\label{ux434ux43eux441ux442ux430ux442ux43eux447ux43dux44bux435-ux443ux441ux43bux43eux432ux438ux44f-ux44dux440ux433ux43eux434ux438ux447ux43dux43eux441ux442ux438}

\textbf{Первое достаточное условие:} корреляции затухают:

\[
R(k)\to0,
\qquad
|k|\to\infty.
\]

Тогда взвешенное среднее ковариаций стремится к нулю, и

\[
\overline X_N\xrightarrow{L^2}m.
\]

\textbf{Второе достаточное условие:} абсолютная суммируемость
ковариаций:

\[
\sum_{k\in\mathbb Z}|R(k)|<\infty.
\]

Тогда

\[
\operatorname{Var}(\overline X_N)
\le
\frac1N\sum_{|k|<N}|R(k)|
\le
\frac1N\sum_{k\in\mathbb Z}|R(k)|
\to0.
\]

Это очень удобное условие: если корреляции убывают достаточно быстро,
среднее эргодично.

\textbf{Третье достаточное условие:} среднее абсолютных корреляций
уходит в ноль:

\[
\frac1N\sum_{k=1}^{N-1}|R(k)|\to0.
\]

Тогда из формулы для дисперсии среднего также следует

\[
\operatorname{Var}(\overline X_N)\to0.
\]

\subsubsection{2.9. Интуитивное понимание и границы
применимости}\label{ux438ux43dux442ux443ux438ux442ux438ux432ux43dux43eux435-ux43fux43eux43dux438ux43cux430ux43dux438ux435-ux438-ux433ux440ux430ux43dux438ux446ux44b-ux43fux440ux438ux43cux435ux43dux438ux43cux43eux441ux442ux438}

Интуитивно все эти условия говорят одно: далекие значения процесса почти
не связаны линейно, поэтому усреднение по времени уничтожает случайные
колебания.

Важно понимать ограничения. Эргодичность по среднему в широком смысле -
это не полная эргодичность в узком смысле. Она касается только средних
самих \(X_n\) и только сходимости в \(L^2\). Узкая эргодичность сильнее
и относится ко всем сдвиг-инвариантным событиям или временным средним
функций \(f(X_n)\).

\subsection{3. Основные
определения}\label{ux43eux441ux43dux43eux432ux43dux44bux435-ux43eux43fux440ux435ux434ux435ux43bux435ux43dux438ux44f}

\subsubsection{Стационарность в широком
смысле}\label{ux441ux442ux430ux446ux438ux43eux43dux430ux440ux43dux43eux441ux442ux44c-ux432-ux448ux438ux440ux43eux43aux43eux43c-ux441ux43cux44bux441ux43bux435}

Последовательность \((X_n)_{n\in\mathbb Z}\) стационарна в широком
смысле, если

\[
\mathbb E|X_n|^2<\infty,
\]

\[
\mathbb EX_n=m
\]

не зависит от \(n\), и

\[
\operatorname{Cov}(X_n,X_m)
\]

зависит только от \(n-m\).

\subsubsection{Ковариационная
функция}\label{ux43aux43eux432ux430ux440ux438ux430ux446ux438ux43eux43dux43dux430ux44f-ux444ux443ux43dux43aux446ux438ux44f}

Для стационарной в широком смысле последовательности:

\[
R(k)=\operatorname{Cov}(X_k,X_0).
\]

В общем нестационарном случае:

\[
R(n,m)=\operatorname{Cov}(X_n,X_m).
\]

\subsubsection{Корреляционная
функция}\label{ux43aux43eux440ux440ux435ux43bux44fux446ux438ux43eux43dux43dux430ux44f-ux444ux443ux43dux43aux446ux438ux44f}

Если \(R(0)>0\), то

\[
\rho(k)=\frac{R(k)}{R(0)}.
\]

Она нормирована:

\[
\rho(0)=1,
\qquad
|\rho(k)|\le1.
\]

\subsubsection{Центрированная
последовательность}\label{ux446ux435ux43dux442ux440ux438ux440ux43eux432ux430ux43dux43dux430ux44f-ux43fux43eux441ux43bux435ux434ux43eux432ux430ux442ux435ux43bux44cux43dux43eux441ux442ux44c}

\[
Y_n=X_n-\mathbb EX_n.
\]

Для стационарной последовательности:

\[
Y_n=X_n-m.
\]

\subsubsection{Эргодичность по среднему в широком
смысле}\label{ux44dux440ux433ux43eux434ux438ux447ux43dux43eux441ux442ux44c-ux43fux43e-ux441ux440ux435ux434ux43dux435ux43cux443-ux432-ux448ux438ux440ux43eux43aux43eux43c-ux441ux43cux44bux441ux43bux435}

Стационарная в широком смысле последовательность эргодична по среднему,
если

\[
\frac1N\sum_{n=0}^{N-1}X_n
\xrightarrow{L^2}
\mathbb EX_0.
\]

\subsubsection{Среднеквадратичная
сходимость}\label{ux441ux440ux435ux434ux43dux435ux43aux432ux430ux434ux440ux430ux442ux438ux447ux43dux430ux44f-ux441ux445ux43eux434ux438ux43cux43eux441ux442ux44c}

\[
Z_N\xrightarrow{L^2}Z
\]

означает

\[
\mathbb E|Z_N-Z|^2\to0.
\]

\subsubsection{Положительная определенность ковариационной
функции}\label{ux43fux43eux43bux43eux436ux438ux442ux435ux43bux44cux43dux430ux44f-ux43eux43fux440ux435ux434ux435ux43bux435ux43dux43dux43eux441ux442ux44c-ux43aux43eux432ux430ux440ux438ux430ux446ux438ux43eux43dux43dux43eux439-ux444ux443ux43dux43aux446ux438ux438}

Функция \(R:\mathbb Z\to\mathbb C\) положительно определена, если

\[
\sum_{\alpha,\beta=1}^r
a_\alpha\overline{a_\beta}R(n_\alpha-n_\beta)\ge0
\]

для любых \(r\), \(a_1,\ldots,a_r\in\mathbb C\) и
\(n_1,\ldots,n_r\in\mathbb Z\).

\subsection{4. Основные теоремы и
свойства}\label{ux43eux441ux43dux43eux432ux43dux44bux435-ux442ux435ux43eux440ux435ux43cux44b-ux438-ux441ux432ux43eux439ux441ux442ux432ux430}

\subsubsection{1. Свойства ковариационной
функции}\label{ux441ux432ux43eux439ux441ux442ux432ux430-ux43aux43eux432ux430ux440ux438ux430ux446ux438ux43eux43dux43dux43eux439-ux444ux443ux43dux43aux446ux438ux438}

Для стационарной в широком смысле последовательности:

\[
R(0)\ge0,
\]

\[
R(-k)=\overline{R(k)},
\]

\[
|R(k)|\le R(0).
\]

Если процесс действительный, то

\[
R(-k)=R(k).
\]

\subsubsection{2. Положительная
определенность}\label{ux43fux43eux43bux43eux436ux438ux442ux435ux43bux44cux43dux430ux44f-ux43eux43fux440ux435ux434ux435ux43bux435ux43dux43dux43eux441ux442ux44c}

Ковариационная функция стационарной последовательности положительно
определена:

\[
\sum_{\alpha,\beta}
a_\alpha\overline{a_\beta}R(n_\alpha-n_\beta)\ge0.
\]

Доказательство:

\[
\sum_{\alpha,\beta}
a_\alpha\overline{a_\beta}R(n_\alpha-n_\beta)
=
\mathbb E\left|
\sum_\alpha a_\alpha(X_{n_\alpha}-m)
\right|^2
\ge0.
\]

\subsubsection{3. Дисперсия временного
среднего}\label{ux434ux438ux441ux43fux435ux440ux441ux438ux44f-ux432ux440ux435ux43cux435ux43dux43dux43eux433ux43e-ux441ux440ux435ux434ux43dux435ux433ux43e}

Пусть

\[
\overline X_N=\frac1N\sum_{n=0}^{N-1}X_n.
\]

Тогда

\[
\operatorname{Var}(\overline X_N)
=
\frac1{N^2}
\sum_{i=0}^{N-1}\sum_{j=0}^{N-1}R(i-j).
\]

Эквивалентно:

\[
\operatorname{Var}(\overline X_N)
=
\frac1N
\sum_{|k|<N}
\left(1-\frac{|k|}{N}\right)R(k).
\]

В действительном случае:

\[
\operatorname{Var}(\overline X_N)
=
\frac{R(0)}{N}
+
\frac{2}{N}\sum_{k=1}^{N-1}
\left(1-\frac{k}{N}\right)R(k).
\]

\subsubsection{4. Критерий эргодичности по
среднему}\label{ux43aux440ux438ux442ux435ux440ux438ux439-ux44dux440ux433ux43eux434ux438ux447ux43dux43eux441ux442ux438-ux43fux43e-ux441ux440ux435ux434ux43dux435ux43cux443}

Последовательность эргодична по среднему в широком смысле тогда и только
тогда, когда

\[
\operatorname{Var}(\overline X_N)\to0.
\]

По формуле выше это равносильно

\[
\frac1{N^2}
\sum_{i,j=0}^{N-1}R(i-j)\to0.
\]

\subsubsection{5. Достаточное условие через затухание
корреляций}\label{ux434ux43eux441ux442ux430ux442ux43eux447ux43dux43eux435-ux443ux441ux43bux43eux432ux438ux435-ux447ux435ux440ux435ux437-ux437ux430ux442ux443ux445ux430ux43dux438ux435-ux43aux43eux440ux440ux435ux43bux44fux446ux438ux439}

Если

\[
R(k)\to0,
\qquad
|k|\to\infty,
\]

то

\[
\overline X_N\xrightarrow{L^2}m.
\]

\subsubsection{6. Достаточное условие через абсолютную
суммируемость}\label{ux434ux43eux441ux442ux430ux442ux43eux447ux43dux43eux435-ux443ux441ux43bux43eux432ux438ux435-ux447ux435ux440ux435ux437-ux430ux431ux441ux43eux43bux44eux442ux43dux443ux44e-ux441ux443ux43cux43cux438ux440ux443ux435ux43cux43eux441ux442ux44c}

Если

\[
\sum_{k\in\mathbb Z}|R(k)|<\infty,
\]

то

\[
\operatorname{Var}(\overline X_N)=O\left(\frac1N\right)
\]

и последовательность эргодична по среднему.

\subsubsection{7. Спектральная
формулировка}\label{ux441ux43fux435ux43aux442ux440ux430ux43bux44cux43dux430ux44f-ux444ux43eux440ux43cux443ux43bux438ux440ux43eux432ux43aux430}

По теореме Герглотца положительно определенная ковариационная функция
имеет представление

\[
R(k)=\int_{-\pi}^{\pi}e^{ik\lambda}F(d\lambda),
\]

где \(F\) - конечная неотрицательная спектральная мера. Тогда средние
сходятся к \(m\) в \(L^2\) тогда и только тогда, когда у спектральной
меры нет атома в нуле:

\[
F(\{0\})=0.
\]

Это уже мост к Билету 25.

\subsection{5. Примеры}\label{ux43fux440ux438ux43cux435ux440ux44b}

\subsubsection{Пример 1. Белый
шум}\label{ux43fux440ux438ux43cux435ux440-1.-ux431ux435ux43bux44bux439-ux448ux443ux43c}

Пусть

\[
\mathbb E X_n=0,
\qquad
\operatorname{Var}X_n=\sigma^2,
\]

и величины некоррелированы:

\[
\operatorname{Cov}(X_n,X_m)=0,
\qquad
n\ne m.
\]

Тогда

\[
R(k)=
\begin{cases}
\sigma^2, & k=0,\\
0, & k\ne0.
\end{cases}
\]

Значит

\[
\operatorname{Var}(\overline X_N)=\frac{\sigma^2}{N}\to0.
\]

Белый шум эргодичен по среднему.

\subsubsection{Пример 2. Геометрически затухающие
корреляции}\label{ux43fux440ux438ux43cux435ux440-2.-ux433ux435ux43eux43cux435ux442ux440ux438ux447ux435ux441ux43aux438-ux437ux430ux442ux443ux445ux430ux44eux449ux438ux435-ux43aux43eux440ux440ux435ux43bux44fux446ux438ux438}

Пусть

\[
R(k)=\sigma^2 a^{|k|},
\qquad
|a|<1.
\]

Тогда

\[
\sum_{k\in\mathbb Z}|R(k)|<\infty.
\]

Следовательно,

\[
\overline X_N\xrightarrow{L^2}m.
\]

Такой вид ковариаций возникает, например, у стационарной авторегрессии
первого порядка \(AR(1)\):

\[
X_n=aX_{n-1}+\varepsilon_n,
\qquad
|a|<1.
\]

\subsubsection{Пример 3. Случайная
константа}\label{ux43fux440ux438ux43cux435ux440-3.-ux441ux43bux443ux447ux430ux439ux43dux430ux44f-ux43aux43eux43dux441ux442ux430ux43dux442ux430}

Пусть

\[
X_n=Z
\]

для всех \(n\), где \(Z\) - нетривиальная случайная величина,

\[
\mathbb EZ=m,
\qquad
\operatorname{Var}Z>0.
\]

Тогда последовательность стационарна, но

\[
\overline X_N=Z
\]

для всех \(N\). Поэтому

\[
\overline X_N\nrightarrow m
\]

в \(L^2\), если \(Z\) не константа. Ковариационная функция:

\[
R(k)=\operatorname{Var}Z
\]

для всех \(k\). Она не затухает, и эргодичности среднего нет.

\subsubsection{Пример 4. Общий случай с суммируемыми
ковариациями}\label{ux43fux440ux438ux43cux435ux440-4.-ux43eux431ux449ux438ux439-ux441ux43bux443ux447ux430ux439-ux441-ux441ux443ux43cux43cux438ux440ux443ux435ux43cux44bux43cux438-ux43aux43eux432ux430ux440ux438ux430ux446ux438ux44fux43cux438}

Если

\[
|R(k)|\le \frac{C}{(1+|k|)^{1+\delta}},
\qquad
\delta>0,
\]

то

\[
\sum_{k\in\mathbb Z}|R(k)|<\infty.
\]

Значит среднее эргодично. Это типичный случай ``быстрого забывания''
прошлого.

\subsubsection{Пример 5. Некоррелированность не равна
независимости}\label{ux43fux440ux438ux43cux435ux440-5.-ux43dux435ux43aux43eux440ux440ux435ux43bux438ux440ux43eux432ux430ux43dux43dux43eux441ux442ux44c-ux43dux435-ux440ux430ux432ux43dux430-ux43dux435ux437ux430ux432ux438ux441ux438ux43cux43eux441ux442ux438}

Можно иметь

\[
R(k)=0
\]

при \(k\ne0\), но величины не будут независимыми. Для эргодичности по
среднему в широком смысле достаточно контроля вторых моментов, но для
узкой эргодичности этого уже мало.

\subsection{6. Схемы доказательств /
обоснований}\label{ux441ux445ux435ux43cux44b-ux434ux43eux43aux430ux437ux430ux442ux435ux43bux44cux441ux442ux432-ux43eux431ux43eux441ux43dux43eux432ux430ux43dux438ux439}

\textbf{Доказательство свойств \(R(k)\).}

Если процесс центрирован, то

\[
R(k)=\mathbb E X_k\overline{X_0}.
\]

Тогда

\[
R(-k)=\mathbb E X_{-k}\overline{X_0}.
\]

По стационарности пара \((X_{-k},X_0)\) имеет те же характеристики
второго порядка, что \((X_0,X_k)\), поэтому

\[
R(-k)=\overline{R(k)}.
\]

Неравенство

\[
|R(k)|\le R(0)
\]

следует из Коши-Буняковского:

\[
|\mathbb E X_k\overline{X_0}|
\le
\sqrt{\mathbb E|X_k|^2\,\mathbb E|X_0|^2}
=R(0).
\]

\textbf{Вывод формулы для дисперсии среднего.}

Центрируем:

\[
Y_n=X_n-m.
\]

Тогда

\[
\overline X_N-m=\frac1N\sum_{n=0}^{N-1}Y_n.
\]

Следовательно,

\[
\mathbb E|\overline X_N-m|^2
=
\frac1{N^2}
\mathbb E
\left|
\sum_{n=0}^{N-1}Y_n
\right|^2.
\]

Раскрываем квадрат:

\[
\mathbb E
\left|
\sum_{n=0}^{N-1}Y_n
\right|^2
=
\sum_{i=0}^{N-1}\sum_{j=0}^{N-1}
\mathbb E Y_i\overline{Y_j}.
\]

По стационарности:

\[
\mathbb E Y_i\overline{Y_j}=R(i-j).
\]

Получаем:

\[
\mathbb E|\overline X_N-m|^2
=
\frac1{N^2}
\sum_{i,j=0}^{N-1}R(i-j).
\]

Если собрать одинаковые лаги:

\[
\mathbb E|\overline X_N-m|^2
=
\frac1N
\sum_{|k|<N}
\left(1-\frac{|k|}{N}\right)R(k).
\]

\textbf{Почему абсолютная суммируемость достаточна.}

Если

\[
\sum_{k\in\mathbb Z}|R(k)|<\infty,
\]

то

\[
\operatorname{Var}(\overline X_N)
\le
\frac1N
\sum_{|k|<N}|R(k)|
\le
\frac1N
\sum_{k\in\mathbb Z}|R(k)|
\to0.
\]

Значит

\[
\overline X_N\to m
\]

в \(L^2\).

\textbf{Почему \(R(k)\to0\) достаточно.}

Из \(|R(k)|\le R(0)\) следует ограниченность ковариаций. Если
\(R(k)\to0\), то средние по Чезаро тоже стремятся к нулю:

\[
\frac1N\sum_{k=0}^{N-1}R(k)\to0.
\]

С учетом симметрии \(R(-k)=\overline{R(k)}\) это дает исчезновение
взвешенного среднего ковариаций в формуле для

\[
\operatorname{Var}(\overline X_N).
\]

Следовательно,

\[
\operatorname{Var}(\overline X_N)\to0.
\]

\subsection{7. Что написать на
доске}\label{ux447ux442ux43e-ux43dux430ux43fux438ux441ux430ux442ux44c-ux43dux430-ux434ux43eux441ux43aux435}

Ковариационная функция:

\[
R(k)=\operatorname{Cov}(X_k,X_0)
=
\mathbb E[(X_k-m)\overline{(X_0-m)}].
\]

Нормированная корреляция:

\[
\rho(k)=\frac{R(k)}{R(0)}.
\]

Свойства:

\[
R(0)\ge0,
\qquad
R(-k)=\overline{R(k)},
\qquad
|R(k)|\le R(0).
\]

Временное среднее:

\[
\overline X_N=\frac1N\sum_{n=0}^{N-1}X_n.
\]

Дисперсия среднего:

\[
\mathbb E|\overline X_N-m|^2
=
\frac1{N^2}
\sum_{i,j=0}^{N-1}R(i-j)
=
\frac1N
\sum_{|k|<N}
\left(1-\frac{|k|}{N}\right)R(k).
\]

Достаточные условия:

\[
R(k)\to0
\quad\Rightarrow\quad
\overline X_N\xrightarrow{L^2}m,
\]

и

\[
\sum_{k\in\mathbb Z}|R(k)|<\infty
\quad\Rightarrow\quad
\overline X_N\xrightarrow{L^2}m.
\]

Спектральная связка:

\[
R(k)=\int_{-\pi}^{\pi}e^{ik\lambda}F(d\lambda),
\qquad
F(\{0\})=0.
\]

\subsection{8. Типичные дополнительные
вопросы}\label{ux442ux438ux43fux438ux447ux43dux44bux435-ux434ux43eux43fux43eux43bux43dux438ux442ux435ux43bux44cux43dux44bux435-ux432ux43eux43fux440ux43eux441ux44b}

\textbf{Чем ковариационная функция отличается от корреляционной?}

Ковариационная функция:

\[
R(k)=\operatorname{Cov}(X_k,X_0).
\]

Нормированная корреляционная функция:

\[
\rho(k)=\frac{R(k)}{R(0)}.
\]

У \(\rho(0)=1\), а у \(R(0)=\operatorname{Var}X_0\).

\textbf{Почему нужно центрирование?}

Корреляция должна измерять совместные отклонения от среднего, поэтому
нужно брать

\[
X_n-m.
\]

Если забыть центрирование, то постоянная ненулевая средняя компонента
даст ложный вклад в зависимость.

\textbf{Почему \(\operatorname{Var}(\overline X_N)\to0\) означает
эргодичность по среднему?}

Потому что

\[
\mathbb E\overline X_N=m.
\]

Значит

\[
\mathbb E|\overline X_N-m|^2
=
\operatorname{Var}(\overline X_N).
\]

Сходимость этой величины к нулю и есть сходимость в \(L^2\).

\textbf{Достаточно ли условия \(R(k)\to0\)?}

Да, для эргодичности среднего в широком смысле это достаточное условие.
Оно не обязательно для всех возможных вариантов, но удобное и часто
используемое.

\textbf{Почему абсолютная суммируемость сильнее?}

Если

\[
\sum_k|R(k)|<\infty,
\]

то ковариации не просто стремятся к нулю, а делают это достаточно
быстро. Тогда дисперсия среднего убывает как \(O(1/N)\).

\textbf{Пример неэргодичной последовательности?}

Случайная константа:

\[
X_n=Z.
\]

Тогда

\[
\overline X_N=Z
\]

для всех \(N\), и среднее не сходится к \(\mathbb EZ\), если \(Z\) не
константа.

\textbf{Связь со спектральной мерой?}

Если

\[
R(k)=\int e^{ik\lambda}F(d\lambda),
\]

то атом спектральной меры в нуле соответствует случайной постоянной
компоненте, которая не исчезает при усреднении. Поэтому условие
эргодичности среднего:

\[
F(\{0\})=0.
\]

\textbf{В чем отличие от узкой эргодичности?}

Узкая эргодичность говорит о сдвиг-инвариантных событиях и временных
средних функций. Эргодичность в широком смысле по среднему говорит
только о сходимости

\[
\frac1N\sum X_n
\]

в \(L^2\).

\subsection{9. Типичные
ошибки}\label{ux442ux438ux43fux438ux447ux43dux44bux435-ux43eux448ux438ux431ux43aux438}

\begin{itemize}
\tightlist
\item
  Путать ковариационную функцию \(R(k)\) и нормированный коэффициент
  корреляции \(\rho(k)\).
\item
  Забывать центрирование: писать \(\mathbb EX_kX_0\) вместо
  \(\mathbb E[(X_k-m)(X_0-m)]\) без условия \(m=0\).
\item
  Говорить ``эргодична'' без уточнения: в узком смысле, по среднему, в
  \(L^2\).
\item
  Считать, что \(R(k)\to0\) означает независимость. Это только
  некоррелированность в пределе.
\item
  Забывать множитель \(N-|k|\) при переходе от двойной суммы к сумме по
  лагам.
\item
  Думать, что стационарность сама по себе дает эргодичность. Случайная
  константа - контрпример.
\item
  Считать абсолютную суммируемость необходимой. Она достаточна, но не
  необходима.
\item
  В комплексном случае забывать сопряжение и свойство
  \(R(-k)=\overline{R(k)}\).
\end{itemize}

\subsection{10. Связи с другими
билетами}\label{ux441ux432ux44fux437ux438-ux441-ux434ux440ux443ux433ux438ux43cux438-ux431ux438ux43bux435ux442ux430ux43cux438}

\begin{itemize}
\tightlist
\item
  Билет 12: ковариация как скалярное произведение в \(L_2\) и
  положительная определенность.
\item
  Билет 14: коэффициент корреляции и геометрия центрированных величин.
\item
  Билет 17: стационарность и эргодичность в широком и узком смыслах.
\item
  Билет 18: для гауссовских последовательностей ковариация определяет
  все конечномерные распределения.
\item
  Билет 21: эргодичность среднего - аналог закона больших чисел для
  зависимых стационарных последовательностей.
\item
  Билет 23: стационарные марковские цепи дают примеры зависимых
  последовательностей, где корреляции часто затухают.
\item
  Билет 25: спектральная мера является преобразованием Фурье
  ковариационной функции.
\end{itemize}

\subsection{11. Минимум на
удовлетворительно}\label{ux43cux438ux43dux438ux43cux443ux43c-ux43dux430-ux443ux434ux43eux432ux43bux435ux442ux432ux43eux440ux438ux442ux435ux43bux44cux43dux43e}

Нужно сказать:

\begin{enumerate}
\def\labelenumi{\arabic{enumi}.}
\tightlist
\item
  Для стационарной в широком смысле последовательности
\end{enumerate}

\[
R(k)=\operatorname{Cov}(X_k,X_0).
\]

\begin{enumerate}
\def\labelenumi{\arabic{enumi}.}
\setcounter{enumi}{1}
\tightlist
\item
  Временное среднее:
\end{enumerate}

\[
\overline X_N=\frac1N\sum_{n=0}^{N-1}X_n.
\]

\begin{enumerate}
\def\labelenumi{\arabic{enumi}.}
\setcounter{enumi}{2}
\tightlist
\item
  Эргодичность по среднему:
\end{enumerate}

\[
\overline X_N\xrightarrow{L^2}\mathbb EX_0.
\]

\begin{enumerate}
\def\labelenumi{\arabic{enumi}.}
\setcounter{enumi}{3}
\tightlist
\item
  Главная формула:
\end{enumerate}

\[
\operatorname{Var}(\overline X_N)
=
\frac1{N^2}\sum_{i,j=0}^{N-1}R(i-j).
\]

\begin{enumerate}
\def\labelenumi{\arabic{enumi}.}
\setcounter{enumi}{4}
\tightlist
\item
  Если корреляции затухают достаточно быстро, например
  \(\sum_k|R(k)|<\infty\), то среднее эргодично.
\end{enumerate}

\subsection{12. Минимум на
хорошо}\label{ux43cux438ux43dux438ux43cux443ux43c-ux43dux430-ux445ux43eux440ux43eux448ux43e}

Помимо минимума, нужно:

\begin{itemize}
\tightlist
\item
  различить \(R(k)\) и \(\rho(k)\);
\item
  назвать свойства \(R(0)\ge0\), \(R(-k)=\overline{R(k)}\),
  \(|R(k)|\le R(0)\);
\item
  вывести формулу
\end{itemize}

\[
\operatorname{Var}(\overline X_N)
=
\frac1N\sum_{|k|<N}\left(1-\frac{|k|}{N}\right)R(k);
\]

\begin{itemize}
\tightlist
\item
  привести пример белого шума и случайной константы;
\item
  объяснить, почему случайная константа стационарна, но не эргодична по
  среднему.
\end{itemize}

\subsection{13. Что добавить для
отлично}\label{ux447ux442ux43e-ux434ux43eux431ux430ux432ux438ux442ux44c-ux434ux43bux44f-ux43eux442ux43bux438ux447ux43dux43e}

Для сильного ответа стоит добавить:

\begin{itemize}
\tightlist
\item
  положительную определенность ковариационной функции;
\item
  условие Чезаро:
\end{itemize}

\[
\frac1N\sum_{k=0}^{N-1}R(k)\to0;
\]

\begin{itemize}
\tightlist
\item
  спектральную интерпретацию:
\end{itemize}

\[
F(\{0\})=0;
\]

\begin{itemize}
\tightlist
\item
  связь с теоремой Герглотца:
\end{itemize}

\[
R(k)=\int_{-\pi}^{\pi}e^{ik\lambda}F(d\lambda);
\]

\begin{itemize}
\tightlist
\item
  четкое различие между эргодичностью в широком смысле и узкой
  эргодичностью.
\end{itemize}

\subsection{14. Устная версия ответа на 5
минут}\label{ux443ux441ux442ux43dux430ux44f-ux432ux435ux440ux441ux438ux44f-ux43eux442ux432ux435ux442ux430-ux43dux430-5-ux43cux438ux43dux443ux442}

Пусть \(X=(X_n)_{n\in\mathbb Z}\) - стационарная в широком смысле
случайная последовательность:

\[
\mathbb EX_n=m,
\qquad
\mathbb E|X_n|^2<\infty,
\]

а ковариация зависит только от лага. Тогда ковариационная функция:

\[
R(k)=\operatorname{Cov}(X_k,X_0)
=
\mathbb E[(X_k-m)\overline{(X_0-m)}].
\]

Если \(R(0)>0\), нормированная корреляционная функция:

\[
\rho(k)=\frac{R(k)}{R(0)}.
\]

Основные свойства:

\[
R(0)\ge0,
\qquad
R(-k)=\overline{R(k)},
\qquad
|R(k)|\le R(0).
\]

Кроме того, \(R\) положительно определена, потому что для любой линейной
комбинации

\[
\mathbb E\left|\sum_\alpha a_\alpha(X_{n_\alpha}-m)\right|^2\ge0.
\]

Эргодичность по среднему в широком смысле означает:

\[
\overline X_N=\frac1N\sum_{n=0}^{N-1}X_n
\xrightarrow{L^2}
m.
\]

Так как \(\mathbb E\overline X_N=m\), это равносильно

\[
\operatorname{Var}(\overline X_N)\to0.
\]

Вычисляем дисперсию:

\[
\operatorname{Var}(\overline X_N)
=
\frac1{N^2}\sum_{i=0}^{N-1}\sum_{j=0}^{N-1}R(i-j).
\]

Собирая одинаковые лаги:

\[
\operatorname{Var}(\overline X_N)
=
\frac1N\sum_{|k|<N}
\left(1-\frac{|k|}{N}\right)R(k).
\]

Отсюда достаточные условия эргодичности: если

\[
R(k)\to0
\]

или тем более

\[
\sum_{k\in\mathbb Z}|R(k)|<\infty,
\]

то

\[
\operatorname{Var}(\overline X_N)\to0,
\]

а значит

\[
\overline X_N\to m
\]

в среднем квадратичном.

Пример эргодичного среднего - белый шум: \(R(k)=0\) при \(k\ne0\), и
дисперсия среднего равна \(\sigma^2/N\). Контрпример - случайная
константа \(X_n=Z\): процесс стационарен, но \(\overline X_N=Z\),
поэтому среднее не сходится к \(\mathbb EZ\), если \(Z\) не вырождена.

В конце можно добавить связь со спектрами: по теореме Герглотца

\[
R(k)=\int_{-\pi}^{\pi}e^{ik\lambda}F(d\lambda).
\]

Эргодичность среднего соответствует отсутствию атома спектральной меры в
нуле:

\[
F(\{0\})=0.
\]

\subsection{15. Если осталось 2 минуты перед
экзаменом}\label{ux435ux441ux43bux438-ux43eux441ux442ux430ux43bux43eux441ux44c-2-ux43cux438ux43dux443ux442ux44b-ux43fux435ux440ux435ux434-ux44dux43aux437ux430ux43cux435ux43dux43eux43c}

Запомнить:

\[
R(k)=\operatorname{Cov}(X_k,X_0).
\]

Свойства:

\[
R(0)\ge0,
\qquad
R(-k)=\overline{R(k)},
\qquad
|R(k)|\le R(0).
\]

Эргодичность среднего:

\[
\frac1N\sum_{n=0}^{N-1}X_n
\xrightarrow{L^2}
\mathbb EX_0.
\]

Главная формула:

\[
\operatorname{Var}(\overline X_N)
=
\frac1{N^2}\sum_{i,j=0}^{N-1}R(i-j).
\]

Достаточные условия:

\[
R(k)\to0
\]

или

\[
\sum_k|R(k)|<\infty.
\]

Главная ловушка: стационарность не означает эргодичность; случайная
константа стационарна, но не эргодична по среднему.

\newpage

\end{document}
